\documentclass{article}
\usepackage{enumitem}
\usepackage{amsmath}
\begin{document}

\title{Chapter 31: Strategies for Enhancement in Food Production}
\date{}
\maketitle

\begin{enumerate}[label=Q\arabic*.]

\item The Green Revolution in India during the 1960s--70s was primarily based on:
\\(A) Organic farming techniques
\\(B) Introduction of high-yielding dwarf varieties of wheat (\textit{Triticum aestivum}) and rice with increased fertilizer and irrigation
\\(C) Genetic engineering of crop plants
\\(D) Polyploidy induction in crop plants

\item Somatic hybridization produces ``somatic hybrids'' by:
\\(A) Crossing two plants sexually
\\(B) Fusing protoplasts of two different plant species using polyethylene glycol or electrical fusion
\\(C) Grafting two plant species together
\\(D) Inducing polyploidy with colchicine

\item The ``Pomato'' (Potato--Tomato hybrid) was produced by:
\\(A) Sexual hybridization
\\(B) Somatic hybridization (protoplast fusion)
\\(C) Genetic engineering
\\(D) Tissue culture

\item Biofortification of crops aims to:
\\(A) Increase crop yield only
\\(B) Improve nutritional quality (vitamins, minerals, proteins, essential fatty acids) of crops through breeding or genetic engineering
\\(C) Make crops resistant to pests
\\(D) Reduce water requirement of crops

\item Assertion: Bt cotton is resistant to bollworm attack.\\
Reason: Bt cotton expresses the Cry (crystal) protein gene from \textit{Bacillus thuringiensis} which is toxic to lepidopteran larvae but not to mammals.
\\(A) Both true, Reason is correct explanation
\\(B) Both true, Reason is not correct explanation
\\(C) Assertion true, Reason false
\\(D) Assertion false, Reason true

\item Totipotency of plant cells is the basis for:
\\(A) Sexual reproduction in plants
\\(B) Micropropagation (tissue culture); a single cell can regenerate into a complete plant
\\(C) Hybridization breeding
\\(D) Mutation breeding

\item Single cell protein (SCP) is produced from:
\\(A) Animal cells in bioreactors
\\(B) Microorganisms (yeast, bacteria, algae) grown on waste materials as a protein-rich food supplement
\\(C) Genetically engineered plant seeds
\\(D) Photosynthetic bacteria only

\item Inbreeding in animals is used to:
\\(A) Increase heterozygosity
\\(B) Fix desirable traits and create pure-breeding lines, but may cause inbreeding depression
\\(C) Maximize hybrid vigor (heterosis)
\\(D) Increase disease resistance

\item MOET (Multiple Ovulation Embryo Transfer Technology) is used in animal husbandry to:
\\(A) Increase milk yield of cows
\\(B) Obtain multiple embryos from genetically superior females by superovulation and transfer to surrogate mothers
\\(C) Produce transgenic animals
\\(D) Treat infertility in cattle

\item Golden Rice is a transgenic variety of rice engineered to produce:
\\(A) Higher protein content
\\(B) $\beta$-carotene (provitamin A) in the endosperm
\\(C) Insect resistance protein (Bt toxin)
\\(D) Drought tolerance proteins

\item Mutation breeding involves:
\\(A) Introducing genes from other organisms
\\(B) Inducing mutations using mutagens (radiation, chemicals) to create new genetic variation in crop plants
\\(C) Crossing two varieties of the same species
\\(D) Tissue culture of existing varieties

\item The high-yielding semi-dwarf wheat variety \textit{Sonalika} was developed in India by:
\\(A) M.S. Swaminathan through hybridization with Norman Borlaug's varieties
\\(B) Genetic engineering
\\(C) Spontaneous mutation
\\(D) Polyploidy

\item Outbreeding in animals refers to:
\\(A) Mating between close relatives
\\(B) Mating of unrelated animals to overcome inbreeding depression and increase heterosis
\\(C) Mating within the same breed
\\(D) Artificial insemination only

\item Match the following improved crop varieties with their characteristics:\\
1. Hybrid maize $\rightarrow$ (a) Heterosis/hybrid vigor for increased yield\\
2. Sharbati Sonora wheat $\rightarrow$ (b) High-yielding semi-dwarf variety from Green Revolution\\
3. Taichung Native 1 rice $\rightarrow$ (c) High-yielding rice variety for tropical conditions\\
4. Golden Rice $\rightarrow$ (d) Transgenic rice with $\beta$-carotene
\\(A) 1-a, 2-b, 3-c, 4-d
\\(B) 1-b, 2-a, 3-d, 4-c
\\(C) 1-c, 2-d, 3-a, 4-b
\\(D) 1-d, 2-c, 3-b, 4-a

\item The term ``heterosis'' (hybrid vigor) refers to:
\\(A) Uniformity in offspring of inbred lines
\\(B) Superior performance of hybrid offspring over both parents, seen in F1 crosses
\\(C) Negative effects of inbreeding
\\(D) Resistance to diseases in polyploid crops

\end{enumerate}
\end{document}
